\documentclass[12pt]{article}
\usepackage[margin=2cm]{geometry}
\setlength{\parindent}{1em}
\setlength{\parskip}{1em}

\usepackage{amsmath}
\usepackage{booktabs,dcolumn,setspace, float}

\begin{document}

%\section{EM algorithm}


\section{A Unified EM Framework for Misclassified Employment Dynamics with Tenure and Nonemployment Durations}

\subsection{States, sampling, and notation}
We observe individuals \(i=1,\dots,N\) at \(t\in\{1,2,3\}\) (quarters). The \emph{true} employment state is \(s_{it}^*\in\{0,1\}\), where 1 denotes employment. The \emph{observed} state is \(s_{it}\in\{0,1\}\), measured with misclassification. Let the latent 3–period history be
\[
H_i=(h_{i1},h_{i2},h_{i3})\in\{0,1\}^3,\qquad h_{it}\equiv s_{it}^*.
\]
We write \(h=(h_1,h_2,h_3)\) when individual indices are inessential, and denote the set of the eight histories by \(\mathcal{H}\).

We also observe \emph{tenure} (employment duration) \(g_{it}\ge 0\) and \emph{nonemployment duration} \(d_{it}\ge 0\). At any \(t\), only one of \((g_{it},d_{it})\) is observed, determined by \(s_{it}\): if \(s_{it}=1\) we observe \(g_{it}\), if \(s_{it}=0\) we observe \(d_{it}\). Survey weights are \(w_i>0\).

\subsection{True dynamics: discrete- and continuous-time views}
\paragraph{Discrete-time Markov chain.} The true chain \(\{s_{t}^*\}\) is first-order Markov with
\[
\alpha\equiv \Pr(s_{1}^*=1),\qquad
\theta_1\equiv \Pr(s_{t}^*=1\mid s_{t-1}^*=1),\qquad
\theta_0\equiv \Pr(s_{t}^*=1\mid s_{t-1}^*=0).
\]
Hence the prior mass on a history \(h\in\mathcal{H}\) is
\begin{equation}
p(H=h\mid\alpha,\theta_0,\theta_1)=
\alpha^{h_1}(1-\alpha)^{1-h_1}\prod_{t=2}^3
\Big[\theta_1^{\mathbf{1}\{h_{t-1}=1,h_t=1\}}(1-\theta_1)^{\mathbf{1}\{h_{t-1}=1,h_t=0\}}
\theta_0^{\mathbf{1}\{h_{t-1}=0,h_t=1\}}(1-\theta_0)^{\mathbf{1}\{h_{t-1}=0,h_t=0\}}\Big].
\label{eq:priorH}
\end{equation}

\paragraph{Continuous-time embedding and duration parameters.}
Assume a two-state continuous-time Markov chain with intensities \(\alpha_1\) (job loss, \(1\to 0\)) and \(\alpha_0\) (job entry, \(0\to 1\)). With quarterly spacing \(\Delta=1\) quarter \(=3\) months, the discrete-time transition probabilities and the continuous-time intensities are linked by
\begin{equation}
\theta_1 = 1 - e^{-\alpha_1\cdot 3},\qquad
\theta_0 = 1 - e^{-\alpha_0\cdot 3}.
\label{eq:ctmc_link}
\end{equation}
The memoryless property implies \emph{true} spell lengths are exponential:
\(
G\sim \mathrm{Exp}(\lambda_g),\ D\sim \mathrm{Exp}(\lambda_d)
\)
with \(\lambda_g=\alpha_1\) and \(\lambda_d=\alpha_0\). Thus, under the \emph{linked} specification,
\begin{equation}
\lambda_g \equiv -\tfrac{1}{3}\log(1-\theta_1),\qquad
\lambda_d \equiv -\tfrac{1}{3}\log(1-\theta_0).
\label{eq:lambdas_link}
\end{equation}
There are two alternatives in estimation: (i) \emph{linked} (enforce \eqref{eq:lambdas_link}), or (ii) \emph{free} (estimate \(\lambda_g,\lambda_d\) directly. We only propose the first, the challenged with the latter are described below.

\subsection{Deterministic clocks for tenure and nonemployment}
Given a history \(h\), the \emph{true} tenure and nonemployment clocks evolve deterministically at quarterly grid points:
\begin{equation}
g_t^*(h)=h_t\big(g_{t-1}^*(h)+\tfrac{1}{4}\big),\quad g_1^*\sim Exp(\lambda_g) ;\qquad
d_t^*(h)=(1-h_t)\big(d_{t-1}^*(h)+\tfrac{1}{4}\big),\quad d_1^*\sim Exp(\lambda_d).
\label{eq:deterministic_clocks}
\end{equation}
Thus a \emph{continuation} \(h_{t-1}=h_t=1\) increments tenure by \(0.25\), and a \emph{start} \(h_{t-1}=0,h_t=1\) resets tenure to \(0.25\). Analogously for nonemployment.

\subsection{Observation model: misclassification and duration measurement}
\paragraph{State misclassification.}
Observed states are misclassified symmetrically with probability \(\pi\):
\begin{equation}
q_{ab}(\pi)\equiv \Pr(s_{t}=a\mid s_{t}^*=b)=
\pi\mathbf{1}\{a\neq b\}+(1-\pi)\mathbf{1}\{a=b\},\qquad a,b\in\{0,1\}.
\label{eq:q_misclass}
\end{equation}

\paragraph{Duration measurement: different variances by regime.}
When \(s_{t}=s_{t}^*\) (correct classification), we observe the \emph{appropriate} clock with additive Gaussian noise; at spell \emph{starts} the marginal law is exGaussian (EMG), and at \emph{continuations} we use Gaussian \emph{increments}. We allow different measurement variances in employment and nonemployment:
\begin{equation}
\varepsilon^g_{t}\sim \mathcal{N}(0,\sigma_g^2),\qquad
\varepsilon^d_{t}\sim \mathcal{N}(0,\sigma_d^2).
\label{eq:regime_sigmas}
\end{equation}

\textbf{Matched employment (\(s_{t}=h_t=1\)).}

\begin{itemize}\itemsep4pt

%\item Start (\(h_{t-1}=0,h_t=1\)): \(g_{it}\) follows the EMG density,
%\[
%g_{it} = X_{it} + \varepsilon^g_{it},\qquad X_{it}\sim \mathrm{Exp}(\lambda_g),
%\]
%with pdf \(f_{\mathrm{EMG}}(g_{it};\lambda_g,\sigma_g^2)\).
\item Start (\(h_{t-1}=0,h_t=1\)): \(g_{t}\) assume the full quarter, with normal noise component,
\[
g_{t} -0.25 = \varepsilon^g_{t},
\]
with pdf \(N(0.25;\sigma_g^2)\).



\item Continuation (\(h_{t-1}=h_t=1\)): the increment
\begin{equation}
\Delta^g_{t} \equiv (g_{t}-g_{t-1}) - \tfrac{1}{4} \ \sim\ \mathcal{N}\!\big(0,\,2\sigma_g^2\big),
\label{eq:emp_increment}
\end{equation}
independently over \(t\) conditional on the history.

\item Employed from the start, (\(h_{1}=1\)): \(g_{t}\) follows the EMG density,
\[
g_{t} = X_{t} + \varepsilon^g_{t},\qquad X_{t}\sim \mathrm{Exp}(\lambda_g )
\]
\end{itemize}




\textbf{Matched nonemployment (\(s_{t}=h_t=0\)).}

Starts and continuations use

\begin{equation}
\Delta^d_{t} \equiv (d_{t}-d_{t-1}) - \tfrac{1}{4} \ \sim\ \mathcal{N}\!\big(0,\,2\sigma_d^2\big).
\label{eq:nonemp_increment}
\end{equation}

While initially nonemployed are EMG with \((\lambda_d,\sigma_d^2)\)

\textbf{Misclassified durations (\(s_{t}\neq h_t\)).}

If \(h_t=1,s_{t}=0\) (true employment, observed nonemployment), the \emph{inappropriate} duration \(D\sim\mathrm{EMG}(\lambda_d, \sigma_d^2)\). If \(h_t=0,s_{t}=1\), then \(G\sim\mathrm{EMG}(\lambda_g, \sigma_g^2)\).

\subsection{Observed-data likelihood}
Let \(y_i=(s_{i1},s_{i2},s_{i3},g_{i1},g_{i2},g_{i3},d_{i1},d_{i2},d_{i3})\). Conditional on \(h\),
\[
p(y_i\mid h;\Theta)=\underbrace{\prod_{t=1}^3 q_{s_{it},h_t}(\pi)}_{\text{state misclassification}}\times
\underbrace{\prod_{t=1}^3 \ell_t(y_{it}\mid s_{it},h_{t-1},h_t;\lambda_g,\lambda_d,\sigma_g^2,\sigma_d^2)}_{\text{duration emission}},
\]
where \(\ell_t\) selects the EMG/Normal-increment/Exponential factor appropriate to \((s_{it},h_{t-1},h_t)\) using \eqref{eq:emp_increment}–\eqref{eq:nonemp_increment}. The prior over histories is \eqref{eq:priorH}. With survey weights \(\{w_i\}\), the observed-data log-likelihood is
\begin{equation}
\ell(\Theta)=\sum_{i=1}^N w_i\log\bigg[\sum_{h\in\mathcal{H}} p(H_i=h\mid\alpha,\theta_0,\theta_1)\,
\prod_{t=1}^3 q_{s_{it},h_t}(\pi)\,
\prod_{t=1}^3 \ell_t(\cdot)\bigg],
\label{eq:obs_ll}
\end{equation}
with parameters
\[
\Theta=
\begin{cases}
(\alpha,\theta_0,\theta_1,\pi,\sigma_g^2,\sigma_d^2) & \text{(linked: }\lambda\text{s via \eqref{eq:lambdas_link})},\\
(\alpha,\theta_0,\theta_1,\pi,\sigma_g^2,\sigma_d^2,\lambda_g,\lambda_d) & \text{(free).}
\end{cases}
\]

Note, this is the standard mixture likelihood for the mixture log-likelihood. 

We assume the following conditional independence:

\begin{itemize}
    \item State misclassifications are independent across periods $t$ \\
    \item Duration emissions are independent across periods $t$ \\
    \item State misclassifications and duration emissions (i.e. observed data generated/emitted by latent state) are independent
\end{itemize}

\subsection{Complete-data formulation}
The observed data likelihood is difficult to optimize due to the sum over latent histories. We follow the the EM-algorithm approach. The complete-data likelihood formulation assumes we know which history, $h$, generated each observation.

Therefore, introduce indicators \(z_{ih}=\mathbf{1}\{H_i=h\}\) with \(\sum_h z_{ih}=1\). Then we have 
\[L_c(\Theta) = \prod_{i=1}^N \prod_{h \in \mathcal{H}} \left[ p(H_i=h\mid\alpha,\theta_0,\theta_1) \prod_{t=1}^3 q_{s_{it},h_t}(\pi) \prod_{t=1}^3 \ell_t(\cdot) \right]^{w_i z_{ih}}\]

Taking the log gives the complete-data log-likelihood as:
\begin{align}
\log L_c(\Theta)
&=\sum_{i=1}^N\sum_{h\in\mathcal{H}} w_i z_{ih}\bigg\{\log p(H_i=h\mid\alpha,\theta_0,\theta_1)
+\sum_{t=1}^3\log q_{s_{it},h_t}(\pi)+\sum_{t=1}^3\log \ell_t(\cdot)\bigg\}.
\label{eq:complete_ll}
\end{align}

The three terms in \eqref{eq:complete_ll} correspond to (i) the Markov prior over histories, (ii) state misclassification, and (iii) duration emissions. The misclassification block depends only on \(\pi\), and \(\alpha\) appears only in the prior block, so both can be maximised independently. However, under the linked specification \eqref{eq:lambdas_link}, the transition probabilities \(\theta_1, \theta_0\) appear in \emph{both} the prior block (through the Markov transition terms) and the emission block (through \(\lambda_g(\theta_1), \lambda_d(\theta_0)\) in the EMG densities). Therefore these two blocks cannot be maximised independently with respect to \(\theta\); the M-step must account for both contributions simultaneously.

\subsection{EM algorithm}
Starting from \(\Theta^{(0)}\), iterate E- and M-steps until convergence of \(\ell(\Theta)\).

%\paragraph{E-step.}
%Compute responsibilities (posterior probabilities of histories):
%\begin{equation}
%\gamma_{ih}\equiv \mathbb{E}[z_{ih}\mid y_i;\Theta^{\text{old}}]
%=\frac{p(H_i=h\mid\alpha,\theta_0,\theta_1)\prod_{t} q_{s_{it},h_t}(\pi)\prod_{t}\ell_t(\cdot)}
%{\sum_{h'\in\mathcal{H}} p(H_i=h'\mid\alpha,\theta_0,\theta_1)\prod_{t} q_{s_{it},h'_t}(\pi)\prod_{t}\ell_t'(\cdot)}.
%\label{eq:resp}
%\end{equation}
%Accumulate weighted sufficient statistics (with \(w_i\gamma_{ih}\) as weights):
%\begin{align}
%C_1 &= \sum_{i,h} w_i\gamma_{ih}\,h_1,\qquad C_0=\sum_{i,h} w_i\gamma_{ih}(1-h_1),\nonumber\\
%D_1 &= \sum_{i,h} w_i\gamma_{ih}\sum_{t=2}^3 \mathbf{1}\{h_{t-1}=1\},\quad
%D_0 = \sum_{i,h} w_i\gamma_{ih}\sum_{t=2}^3 \mathbf{1}\{h_{t-1}=0\},\nonumber\\
%T_{11} &= \sum_{i,h} w_i\gamma_{ih}\sum_{t=2}^3 \mathbf{1}\{h_{t-1}=1,h_t=1\},\quad
%T_{01} = \sum_{i,h} w_i\gamma_{ih}\sum_{t=2}^3 \mathbf{1}\{h_{t-1}=0,h_t=1\},\nonumber\\
%M &= \sum_{i,h} w_i\gamma_{ih}\sum_{t=1}^3 \mathbf{1}\{s_{it}\neq h_t\}. \label{eq:suff_trans}
%\end{align}
%For the duration block, separate matched continuations, matched starts, and misclassified:
%\begin{align}
%S_g &= \sum_{i,h} w_i\gamma_{ih}\sum_{t:\,h_{t-1}=h_t=1=s_{it}} \Big(\Delta^g_{it}\Big)^2,\quad
%N_g = \sum_{i,h} w_i\gamma_{ih}\sum_{t:\,h_{t-1}=h_t=1=s_{it}} 1,\nonumber\\
%S_d &= \sum_{i,h} w_i\gamma_{ih}\sum_{t:\,h_{t-1}=h_t=0=s_{it}} \Big(\Delta^d_{it}\Big)^2,\quad
%N_d = \sum_{i,h} w_i\gamma_{ih}\sum_{t:\,h_{t-1}=h_t=0=s_{it}} 1, \label{eq:suff_gauss}
%\end{align}
%with \(\Delta^g_{it},\Delta^d_{it}\) from \eqref{eq:emp_increment}–\eqref{eq:nonemp_increment}. For misclassification exponentials,
%\begin{align}
%K_g &= \sum_{i,h} w_i\gamma_{ih}\sum_{t:\,h_t=0,\,s_{it}=1} g_{it},\qquad
%n_g = \sum_{i,h} w_i\gamma_{ih}\sum_{t:\,h_t=0,\,s_{it}=1} 1,\nonumber\\
%K_d &= \sum_{i,h} w_i\gamma_{ih}\sum_{t:\,h_t=1,\,s_{it}=0} d_{it},\qquad
%n_d = \sum_{i,h} w_i\gamma_{ih}\sum_{t:\,h_t=1,\,s_{it}=0} 1. \label{eq:suff_exp}
%\end{align}
%Finally, collect matched \emph{starts} (for EMG), e.g.
%\begin{align}
%\mathcal{S}_g &= \{(i,t,h): h_{t-1}=0,h_t=1,\,s_{it}=1\},\qquad
%\mathcal{S}_d = \{(i,t,h): h_{t-1}=1,h_t=0,\,s_{it}=0\}, \label{eq:start_sets}
%\end{align}
%and the corresponding weighted samples \(g_{it}\) or \(d_{it}\) with weights \(w_i\gamma_{ih}\).


%-------------------------------------



\paragraph{E-step.}

Starting from \(\Theta^{(0)}\), iterate E- and M-steps until convergence of \(\ell(\Theta)\).

Compute responsibilities (posterior probabilities of histories):
\begin{equation}
\gamma_{ih}\equiv \mathbb{E}[z_{ih}\mid y_i;\Theta^{\text{old}}]
=\frac{p(H_i=h\mid\alpha^{\text{old}},\theta_0^{\text{old}},\theta_1^{\text{old}})\prod_{t=1}^3 q_{s_{it},h_t}(\pi^{\text{old}})\prod_{t=1}^3\ell_t(y_{it}\mid s_{it},h_{t-1},h_t;\Theta^{\text{old}})}
{\sum_{h'\in\mathcal{H}} p(H_i=h'\mid\alpha^{\text{old}},\theta_0^{\text{old}},\theta_1^{\text{old}})\prod_{t=1}^3 q_{s_{it},h'_t}(\pi)\prod_{t=1}^3\ell_t(y_{it}\mid s_{it},h'_{t-1},h'_t;\Theta^{\text{old}})},
\label{eq:resp}
\end{equation}
where \(\ell_t(y_{it}\mid s_{it},h_{t-1},h_t;\Theta)\) is the duration emission density that depends on the observation type:

\begin{itemize}
\item \textbf{Wave~1 matched employment} (\(t=1\), \(h_1=1\), \(s_{1}=1\)):
\[
\ell_1(\cdot) = f_{\mathrm{EMG}}(g_{1};\lambda_g,\sigma_g^2),
\]


\item \textbf{Wave~1 matched nonemployment} (\(t=1\), \(h_1=0\), \(s_{1}=0\)):
\[
\ell_1(\cdot) = f_{\mathrm{EMG}}(d_{1};\lambda_d,\sigma_d^2).
\]
\item \textbf{Employment continuation, previous observed} (\(t\ge 2\), \(h_{t-1}=1\), \(h_t=1\), \(s_{it}=1\), \(s_{i,t-1}=1\)):
\[
\ell_t(\cdot) = \phi\!\left(\frac{\Delta^g_{t}}{\sqrt{2}\sigma_g}\right)\Big/\!\big(\sqrt{2}\sigma_g\big),
\quad\text{where}\quad \Delta^g_{t} = (g_{t}-g_{t-1}) - 0.25.
\]

\item \textbf{Employment continuation, previous misclassified} (\(t\ge 2\), \(h_{t-1}=1\), \(h_t=1\), \(s_{t}=1\), \(s_{t-1}=0\)):
\[
\ell_t(\cdot) = f_{\mathrm{EMG}}(g_{t};\lambda_g,\sigma_g^2).
\]

\item \textbf{Nonemployment continuation, previous observed} (\(t\ge 2\), \(h_{t-1}=0\), \(h_t=0\), \(s_{t}=0\), \(s_{t-1}=0\)):
\[
\ell_t(\cdot) = \phi\!\left(\frac{\Delta^d_{t}}{\sqrt{2}\sigma_d}\right)\Big/\!\big(\sqrt{2}\sigma_d\big),
\quad\text{where}\quad \Delta^d_{t} = (d_{t}-d_{t-1}) - 0.25.
\]

\item \textbf{Nonemployment continuation, previous misclassified} (\(t\ge 2\), \(h_{t-1}=0\), \(h_t=0\), \(s_{t}=0\), \(s_{t-1}=1\)):
\[
\ell_t(\cdot) = f_{\mathrm{EMG}}(d_{t};\lambda_d,\sigma_d^2).
\]


\item \textbf{Within-panel employment start} (\(t\ge 2\), \(h_{t-1}=0\), \(h_t=1\), \(s_{it}=1\)):
\[
\ell_t(\cdot) = \phi\!\left(\frac{g_{t}-0.25}{\sigma_g}\right)\Big/\sigma_g.
\]

\item \textbf{Within-panel nonemployment start} (\(t\ge 2\), \(h_{t-1}=1\), \(h_t=0\), \(s_{t}=0\)):
\[
\ell_t(\cdot) = \phi\!\left(\frac{d_{t}-0.25}{\sigma_d}\right)\Big/\sigma_d.
\]

\item \textbf{Misclassified as employed} (\(h_t=0\), \(s_{t}=1\)):
\[
\ell_t(\cdot) = f_{\mathrm{EMG}}(g_{t};\lambda_g,\sigma_g^2).
\]

\item \textbf{Misclassified as nonemployed} (\(h_t=1\), \(s_{t}=0\)):
\[
\ell_t(\cdot) = f_{\mathrm{EMG}}(d_{t};\lambda_d,\sigma_d^2).
\]
\end{itemize}

Under the \textbf{linked specification}, \(\lambda_g\) and \(\lambda_d\) are computed from the transition probabilities via \eqref{eq:lambdas_link} at each iteration.


\medskip

Accumulate weighted sufficient statistics using \(w_i\gamma_{ih}\) as weights:

\textbf{Markov transition statistics:}
\begin{align}
C_1 &= \sum_{i,h} w_i\gamma_{ih}\,h_1,\qquad C_0=\sum_{i,h} w_i\gamma_{ih}(1-h_1),\nonumber\\
D_1 &= \sum_{i,h} w_i\gamma_{ih}\sum_{t=2}^3 \mathbf{1}\{h_{t-1}=1\},\quad
D_0 = \sum_{i,h} w_i\gamma_{ih}\sum_{t=2}^3 \mathbf{1}\{h_{t-1}=0\},\nonumber\\
T_{11} &= \sum_{i,h} w_i\gamma_{ih}\sum_{t=2}^3 \mathbf{1}\{h_{t-1}=1,h_t=1\},\quad
T_{01} = \sum_{i,h} w_i\gamma_{ih}\sum_{t=2}^3 \mathbf{1}\{h_{t-1}=0,h_t=1\},\nonumber\\
M &= \sum_{i,h} w_i\gamma_{ih}\sum_{t=1}^3 \mathbf{1}\{s_{it}\neq h_t\}. \label{eq:suff_trans}
\end{align}

\textbf{Duration variance statistics:}

Two sources of information identify the measurement variances \(\sigma_g^2\) and \(\sigma_d^2\): continuation increments (where both current and lagged durations are observed) and within-panel starts (where the level is observed with known mean).

\emph{Employment continuations} (\(t\ge 2\), \(h_{t-1}=h_t=1\), \(s_{it}=1\), \(s_{i,t-1}=1\)):
\begin{align}
S_g &= \sum_{i,h} w_i\gamma_{ih}\sum_{t=2}^3\mathbf{1}\{h_{t-1}=h_t=1,\,s_{it}=1,\,s_{i,t-1}=1\} \big(\Delta^g_{it}\big)^2,\nonumber\\
N_g &= \sum_{i,h} w_i\gamma_{ih}\sum_{t=2}^3\mathbf{1}\{h_{t-1}=h_t=1,\,s_{it}=1,\,s_{i,t-1}=1\},
\label{eq:suff_gauss_g}
\end{align}
where \(\Delta^g_{it} = (g_{it}-g_{i,t-1}) - 0.25\) from \eqref{eq:emp_increment}. The condition \(s_{i,t-1}=1\) is required so that the lagged tenure \(g_{i,t-1}\) is observed. When the previous wave is misclassified (\(s_{i,t-1}=0\)), the increment is undefined and that observation is excluded. The efficiency loss is proportional to \(\pi\) and small. The estimator remains unbiased because \(\sigma_g^2\) is independent of classification at \(t-1\).

\emph{Nonemployment continuations} (\(t\ge 2\), \(h_{t-1}=h_t=0\), \(s_{it}=0\), \(s_{i,t-1}=0\)):
\begin{align}
S_d &= \sum_{i,h} w_i\gamma_{ih}\sum_{t=2}^3\mathbf{1}\{h_{t-1}=h_t=0,\,s_{it}=0,\,s_{i,t-1}=0\} \big(\Delta^d_{it}\big)^2,\nonumber\\
N_d &= \sum_{i,h} w_i\gamma_{ih}\sum_{t=2}^3\mathbf{1}\{h_{t-1}=h_t=0,\,s_{it}=0,\,s_{i,t-1}=0\},
\label{eq:suff_gauss_d}
\end{align}
where \(\Delta^d_{it} = (d_{it}-d_{i,t-1}) - 0.25\) from \eqref{eq:nonemp_increment}. The same logic applies: \(s_{i,t-1}=0\) ensures the lagged nonemployment duration is observed.

\emph{Within-panel employment starts} (\(t\ge 2\), \(h_{t-1}=0\), \(h_t=1\), \(s_{it}=1\)):
\begin{align}
S_g^{\mathrm{start}} &= \sum_{i,h} w_i\gamma_{ih}\sum_{t=2}^3\mathbf{1}\{h_{t-1}=0,\,h_t=1,\,s_{it}=1\}\big(g_{it}-0.25\big)^2,\nonumber\\
N_g^{\mathrm{start}} &= \sum_{i,h} w_i\gamma_{ih}\sum_{t=2}^3\mathbf{1}\{h_{t-1}=0,\,h_t=1,\,s_{it}=1\}.
\label{eq:suff_gauss_g_start}
\end{align}

\emph{Within-panel nonemployment starts} (\(t\ge 2\), \(h_{t-1}=1\), \(h_t=0\), \(s_{it}=0\)):
\begin{align}
S_d^{\mathrm{start}} &= \sum_{i,h} w_i\gamma_{ih}\sum_{t=2}^3\mathbf{1}\{h_{t-1}=1,\,h_t=0,\,s_{it}=0\}\big(d_{it}-0.25\big)^2,\nonumber\\
N_d^{\mathrm{start}} &= \sum_{i,h} w_i\gamma_{ih}\sum_{t=2}^3\mathbf{1}\{h_{t-1}=1,\,h_t=0,\,s_{it}=0\}.
\label{eq:suff_gauss_d_start}
\end{align}

\textbf{Remark on within-panel starts.} Within-panel starts follow \(N(0.25, \sigma^2)\) and yield closed-form sufficient statistics \((S_g^{\mathrm{start}}, N_g^{\mathrm{start}}, S_d^{\mathrm{start}}, N_d^{\mathrm{start}})\), which are pooled with the continuation statistics in the M-step for \(\sigma_g^2\) and \(\sigma_d^2\).

\textbf{EMG gradient statistics for the joint \(\theta\) update.}

Under the linked specification \eqref{eq:lambdas_link}, the EMG emission terms depend on \(\theta_1\) through \(\lambda_g(\theta_1) = -\tfrac{1}{3}\log(1-\theta_1)\), and on \(\theta_0\) through \(\lambda_d(\theta_0) = -\tfrac{1}{3}\log(1-\theta_0)\). Since these terms appear in the Q function alongside the Markov prior, the M-step for \(\theta\) must account for both contributions. The EMG log-density is
\begin{equation}
\log f_{\mathrm{EMG}}(x; \lambda, \sigma^2) = \log \lambda + \tfrac{\lambda^2 \sigma^2}{2} - \lambda x + \log \Phi\!\left(\frac{x - \lambda\sigma^2}{\sigma}\right) + \mathrm{const},
\label{eq:log_emg}
\end{equation}
where \(\Phi(\cdot)\) is the standard Normal CDF. Its derivative with respect to \(\lambda\) is
\begin{equation}
\frac{\partial \log f_{\mathrm{EMG}}}{\partial \lambda} = \frac{1}{\lambda} + \lambda\sigma^2 - x - \sigma\,\frac{\phi(u)}{\Phi(u)},
\label{eq:demg_dlambda}
\end{equation}
where \(u = (x - \lambda\sigma^2)/\sigma\) and \(\phi(u)/\Phi(u)\) is the inverse Mills ratio.

Six observation types contribute EMG terms involving \(\lambda_g\) (and symmetrically for \(\lambda_d\)):
\begin{enumerate}
\item \textbf{Wave~1 matched employment} (\(t=1\), \(h_1=1\), \(s_{i1}=1\)): contributes \(\log f_{\mathrm{EMG}}(g_{i1}; \lambda_g, \sigma_g^2)\).
\item \textbf{Employment continuation, previous misclassified} (\(t\ge 2\), \(h_{t-1}=1\), \(h_t=1\), \(s_{it}=1\), \(s_{i,t-1}=0\)): contributes \(\log f_{\mathrm{EMG}}(g_{it}; \lambda_g, \sigma_g^2)\).
\item \textbf{Misclassified as employed} (\(h_t=0\), \(s_{it}=1\)): contributes \(\log f_{\mathrm{EMG}}(g_{it}; \lambda_g, \sigma_g^2)\).
\end{enumerate}
The corresponding three types for \(\lambda_d\) are: wave~1 matched nonemployment (\(h_1=0\), \(s_{i1}=0\)); nonemployment continuation, previous misclassified (\(h_{t-1}=0\), \(h_t=0\), \(s_{it}=0\), \(s_{i,t-1}=1\)); and misclassified as nonemployed (\(h_t=1\), \(s_{it}=0\)).

Define the weighted EMG gradient contributions accumulated during the E-step:
\begin{align}
E_g &= \sum_{(i,t,h) \in \mathcal{E}_g} w_i\gamma_{ih}\,\frac{\partial \log f_{\mathrm{EMG}}(x_{it}; \lambda_g, \sigma_g^2)}{\partial \lambda_g}, \nonumber\\
E_d &= \sum_{(i,t,h) \in \mathcal{E}_d} w_i\gamma_{ih}\,\frac{\partial \log f_{\mathrm{EMG}}(x_{it}; \lambda_d, \sigma_d^2)}{\partial \lambda_d},
\label{eq:suff_emg_grad}
\end{align}
where \(\mathcal{E}_g\) (resp.\ \(\mathcal{E}_d\)) denotes the set of all \((i,t,h)\) triples belonging to the three observation types above, and \(x_{it}\) is the observed duration (\(g_{it}\) or \(d_{it}\)) at that wave. The partial derivatives are evaluated at the \emph{current} parameter values \(\lambda_g^{\text{old}}, \sigma_g^{2,\text{old}}\) (resp.\ \(\lambda_d^{\text{old}}, \sigma_d^{2,\text{old}}\)).



\paragraph{Free specification (which we don't use).}

If one were to estimate \(\lambda_g, \lambda_d\) independently of the transition probabilities (``free mode''), the M-step would require numerical optimization because all duration observations involving \(\lambda\) use the EMG distribution, which has no closed-form MLE. Specifically, one would need to maximize:
\begin{align}
Q(\lambda_g) = \sum_{i,h} w_i\gamma_{ih}\bigg[&\mathbf{1}\{h_1=1,\,s_{i1}=1\}\log f_{\mathrm{EMG}}(g_{i1};\lambda_g,\sigma_g^2) \nonumber\\
+\sum_{t=2}^3 &\mathbf{1}\{h_{t-1}=1,\,h_t=1,\,s_{it}=1,\,s_{i,t-1}=0\}\log f_{\mathrm{EMG}}(g_{it};\lambda_g,\sigma_g^2) \nonumber\\
+\sum_{t=1}^3 &\mathbf{1}\{h_t=0,\,s_{it}=1\}\log f_{\mathrm{EMG}}(g_{it};\lambda_g,\sigma_g^2)\bigg]
\label{eq:Q_lambda_free}
\end{align}

using numerical methods at each M-step. This approach is:
\begin{itemize}
\item Computationally expensive (nested optimization within EM)
\item Potentially poorly identified (EMG confounds \(\lambda\) and \(\sigma^2\))
\item Theoretically redundant (the CTMC link provides strong structural information)
\end{itemize}

We instead focus on the theoretically justified linked specification.










%-----------------------------------










\paragraph{M-step.}
The E-step allows us to compute the responsibilities, $\gamma_{ih}$, as in eq \eqref{eq:resp}, which is the posterior probability that an individual $i$
 has the true history $h$ given the data and current parameters $\Theta^{\text{old}}$. The soft assignments allow us to compute the expected complete data log-likelihood, $Q(\Theta\mid\Theta^{\text{old}})$, and maximize it treating the responsibilities as fixed.

Maximizing \(Q(\Theta\mid\Theta^{\text{old}})=\mathbb{E}[\log L_c(\Theta)\mid \{y_i\};\Theta^{\text{old}}]\), which amounts to replacing \(z_{ih}\) by \(\gamma_{ih}\) in \eqref{eq:complete_ll}. The Q function decomposes additively into three blocks. The misclassification block (\(\pi\)) and the measurement-variance block (\(\sigma_g^2, \sigma_d^2\)) can each be maximised independently in closed form. However, under the linked specification, \(\theta_1\) and \(\theta_0\) enter both the Markov prior block and the EMG emission block (via \(\lambda_g(\theta_1), \lambda_d(\theta_0)\)), so the M-step for \(\theta\) must jointly account for both contributions. We derive the joint first-order conditions below.

\textbf{Initial-state probability:}
\begin{equation}
\hat{\alpha}=\frac{C_1}{C_0+C_1}.
\label{eq:mstep_alpha}
\end{equation}
If stationarity is imposed, set \(\alpha \leftarrow \theta_0/(\theta_0+1-\theta_1)\).

\textbf{Misclassification:}
\begin{equation}
\hat{\pi}=\frac{M}{3\sum_i w_i}.
\label{eq:mstep_pi}
\end{equation}



\textbf{Measurement variances (pooling continuations and starts):}

Continuations contribute \(\Delta^g_{it}\sim N(0,2\sigma_g^2)\) and starts contribute \((g_{it}-0.25)\sim N(0,\sigma_g^2)\). Writing \(v=\sigma_g^2\), the joint profile log-likelihood is
\[
\ell(v) = -\tfrac{N_g}{2}\log(2v) - \frac{S_g}{4v} - \tfrac{N_g^{\mathrm{start}}}{2}\log(v) - \frac{S_g^{\mathrm{start}}}{2v}.
\]
Setting \(\partial\ell/\partial v=0\) and multiplying through by \(2v^2\) gives
\[
(N_g+N_g^{\mathrm{start}})\,v = \tfrac{1}{2}S_g + S_g^{\mathrm{start}},
\]
so the pooled MLE is
\begin{equation}
\widehat{\sigma}_g^2=\frac{S_g + 2\,S_g^{\mathrm{start}}}{2(N_g + N_g^{\mathrm{start}})},\qquad
\widehat{\sigma}_d^2=\frac{S_d + 2\,S_d^{\mathrm{start}}}{2(N_d + N_d^{\mathrm{start}})}.
\label{eq:mstep_sigmas}
\end{equation}
The factor of 2 on \(S^{\mathrm{start}}\) rescales start contributions (variance \(\sigma^2\)) to the continuation scale (variance \(2\sigma^2\)). When \(N^{\mathrm{start}}=0\), this reduces to \(\hat{\sigma}_g^2 = S_g/(2N_g)\). Enforce small floors if needed.







%------------


\textbf{Transition probabilities and exponential rates (joint update):}

Under the linked specification, \(\theta_1\) enters both the Markov prior block (via the transition counts \(T_{11}, D_1\)) and the emission block (via \(\lambda_g(\theta_1)\) in the EMG densities). The M-step must therefore maximise \(Q\) jointly with respect to \(\theta_1\), collecting gradient contributions from both blocks. An analogous derivation applies to \(\theta_0\).

Using \(\lambda_g = -\tfrac{1}{3}\log(1-\theta_1)\), we have
\[
\frac{d\lambda_g}{d\theta_1} = \frac{1}{3(1-\theta_1)}.
\]

The first-order condition \(\partial Q / \partial \theta_1 = 0\) collects two terms:
\begin{equation}
\underbrace{\frac{T_{11}}{\theta_1} - \frac{D_1 - T_{11}}{1-\theta_1}}_{\text{from } Q_{\mathrm{prior}}} \;+\; \underbrace{\frac{1}{3(1-\theta_1)} \cdot E_g}_{\text{from } Q_{\mathrm{emission}}} \;=\; 0,
\label{eq:foc_theta1}
\end{equation}
where \(E_g\) is the weighted EMG gradient contribution defined in \eqref{eq:suff_emg_grad}. If the EMG terms were absent (\(E_g = 0\)), this reduces to the standard closed-form update \(\hat{\theta}_1 = T_{11}/D_1\). With \(E_g \neq 0\), the equation has no closed-form solution because \(E_g\) depends on \(\lambda_g(\theta_1)\) through both the exponential and inverse Mills ratio terms in \eqref{eq:demg_dlambda}.

We solve \eqref{eq:foc_theta1} by one-dimensional Newton--Raphson on \(\theta_1 \in (0,1)\), seeded at the prior-only estimate \(T_{11}/D_1\). The second derivative (Hessian) combines the prior curvature
\[
\frac{\partial^2 Q_{\mathrm{prior}}}{\partial \theta_1^2} = -\frac{T_{11}}{\theta_1^2} - \frac{D_1 - T_{11}}{(1-\theta_1)^2}
\]
with the emission curvature obtained by differentiating \eqref{eq:demg_dlambda} again. The Newton step is
\[
\theta_1^{(k+1)} = \theta_1^{(k)} - \left(\frac{\partial^2 Q}{\partial \theta_1^2}\right)^{\!-1} \frac{\partial Q}{\partial \theta_1},
\]
projected onto \((\varepsilon, 1-\varepsilon)\) for numerical safety. In practice, convergence to machine precision requires 2--5 Newton iterations because the seed \(T_{11}/D_1\) is already close to the optimum.

The analogous FOC for \(\theta_0\) is
\begin{equation}
\frac{T_{01}}{\theta_0} - \frac{D_0 - T_{01}}{1-\theta_0} + \frac{1}{3(1-\theta_0)} \cdot E_d = 0,
\label{eq:foc_theta0}
\end{equation}
solved identically by Newton--Raphson seeded at \(T_{01}/D_0\).

Once \(\hat{\theta}_1\) and \(\hat{\theta}_0\) are obtained, the exponential rates follow deterministically:
\begin{equation}
\hat{\lambda}_g = -\tfrac{1}{3}\log(1-\hat{\theta}_1),\qquad
\hat{\lambda}_d = -\tfrac{1}{3}\log(1-\hat{\theta}_0).
\label{eq:mstep_lambda}
\end{equation}

This joint update ensures that the M-step maximises \(Q\) with respect to \(\theta\) (and hence \(\lambda\)) globally, restoring the standard EM monotonicity guarantee.

Under the free specification (not implemented), one would instead maximise \(Q(\lambda_g)\) in \eqref{eq:Q_lambda_free} numerically at each M-step, since the EMG distribution has no closed-form MLE for the rate parameter.





%\textbf{Exponential rates.} Two alternatives:

%\noindent\underline{Linked (CTMC) model:} set \(\lambda_g,\lambda_d\) by \eqref{eq:lambdas_link} using \((\hat\theta_1,\hat\theta_0)\). No separate M-step is needed.

%\noindent\underline{Semi-linked/free model:} use misclassification draws and matched starts.
%\begin{itemize}\itemsep4pt
%\item From misclassification: \(\hat{\lambda}_g^{\text{(mis)}}=n_g/K_g,\ \hat{\lambda}_d^{\text{(mis)}}=n_d/K_d\).
%\item From matched starts (EMG): either (A) \emph{maximize} the EMG contribution numerically in \((\lambda_g,\sigma_g^2)\) and in \((\lambda_d,\sigma_d^2)\) using weighted samples \(\mathcal{S}_g,\mathcal{S}_d\); or (B) \emph{augment} by latent exponentials at starts (call them \(X_{it}\) and \(Y_{it}\)) and set
%\[
%\hat{\lambda}_g^{\text{(start)}}=\frac{\sum_{(i,t,h)\in\mathcal{S}_g} w_i\gamma_{ih}}{\sum_{(i,t,h)\in\mathcal{S}_g} w_i\gamma_{ih}\,\mathbb{E}[X_{it}\mid g_{it},\lambda_g,\sigma_g^2]},\quad
%\hat{\lambda}_d^{\text{(start)}}=\frac{\sum_{(i,t,h)\in\mathcal{S}_d} w_i\gamma_{ih}}{\sum_{(i,t,h)\in\mathcal{S}_d} w_i\gamma_{ih}\,\mathbb{E}[Y_{it}\mid d_{it},\lambda_d,\sigma_d^2]},
%\]
%where the EMG conditional expectations have closed form via Mills ratios.
%\end{itemize}
%Combine the two sources (e.g. inverse-variance weighting or joint maximization of their summed log-likelihoods). If one prefers to softly \emph{link} \(\lambda\)s to \(\theta\)s, add a quadratic penalty \(\rho\big(\lambda_g + \frac{1}{3}\log(1-\theta_1)\big)^2+\rho\big(\lambda_d + \frac{1}{3}\log(1-\theta_0)\big)^2\).



%-----------------------------------




\paragraph{Iteration and convergence.}
Iterate \((\text{E-step},\text{M-step})\) until \(\ell(\Theta)\) in \eqref{eq:obs_ll} increases by less than a tolerance. The E-step uses \eqref{eq:resp}; the M-step applies \eqref{eq:mstep_alpha}, \eqref{eq:mstep_pi}, \eqref{eq:mstep_sigmas}, \eqref{eq:foc_theta1}--\eqref{eq:foc_theta0}, and \eqref{eq:mstep_lambda}. Monotone ascent is guaranteed because every M-step update maximises \(Q\): \(\alpha\), \(\pi\), and \((\sigma_g^2, \sigma_d^2)\) in closed form; \((\theta_1, \theta_0)\) via Newton--Raphson on the joint first-order conditions \eqref{eq:foc_theta1}--\eqref{eq:foc_theta0} that account for both the Markov prior and EMG emission contributions.

\newpage

\subsection{Identification and practical considerations}
\begin{itemize}\itemsep4pt
\item \emph{CTMC link and joint identification.} Under the linked specification, the exponential rates \(\lambda_g,\lambda_d\) are not free parameters but are determined by the transition probabilities \(\theta_1,\theta_0\) via \eqref{eq:lambdas_link}. This ties the discrete-time Markov dynamics and the continuous-time duration model into a single consistent framework, reducing the parameter count and avoiding separate (poorly identified) estimation of the rates.
\item \emph{What identifies \(\pi\)?} Duration data provides strong identification of misclassification. If a person reports employment at all three waves but with tenure increments inconsistent with \(N(0,2\sigma_g^2)\), the model attributes this to misclassification rather than measurement noise. The EMG emission densities for misclassified observations have a different shape from the Normal densities used for correctly classified observations, allowing the two to be distinguished.
\item \emph{What identifies \(\sigma_g^2,\sigma_d^2\)?} Continuation increments from correctly classified consecutive waves identify \(\sigma_g^2\) and \(\sigma_d^2\) directly via \eqref{eq:mstep_sigmas}, independently of \(\pi\) and \(\lambda\). Within-panel starts provide additional information. These two sources are pooled in the M-step.
\item \emph{Overidentification.} With 3 waves there are \(2^3=8\) observed state patterns and continuous duration data at each wave. The linked specification has 6 parameters \((\alpha,\theta_0,\theta_1,\pi,\sigma_g^2,\sigma_d^2)\). The model is overidentified, which allows for specification testing.
\item \emph{Different measurement variances.} Allowing \(\sigma_g^2\neq\sigma_d^2\) captures regime-specific noise and prevents the model from attributing asymmetric measurement error to \(\pi\) or the transition probabilities.
\item \emph{Increment likelihoods.} Using \eqref{eq:emp_increment}--\eqref{eq:nonemp_increment} differences out the unknown true duration level, leaving only measurement noise. This yields a clean \(N(0,2\sigma^2)\) density that identifies \(\sigma^2\) without dependence on \(\lambda\), and avoids using the same measurement error twice in consecutive level likelihoods.
\item \emph{Regularization.} Small floors on \(\sigma_g^2,\sigma_d^2\) guard against pathological solutions typical in mixture-like settings. Transition caps (e.g., \(\theta_1\le 0.995\)) can be imposed if needed.
\item \emph{Weights.} All sufficient statistics are weight-adjusted via \(w_i\gamma_{ih}\); likewise, the observed-data log-likelihood \eqref{eq:obs_ll} uses weights.
\end{itemize}

\newpage



\subsection{Summary of parameterizations}
\begin{center}
\begin{tabular}{lll}
\toprule
Block & Update & Notes \\
\midrule
\(\alpha\) & Closed form \eqref{eq:mstep_alpha} & Stationary option available.\\
\((\theta_1,\theta_0)\) & Newton--Raphson \eqref{eq:foc_theta1}--\eqref{eq:foc_theta0} & Joint FOC: Markov + EMG terms. 2--5 iters.\\
\(\pi\) & Closed form \eqref{eq:mstep_pi} & \\
\((\sigma_g^2,\sigma_d^2)\) & Closed form \eqref{eq:mstep_sigmas} & Pools continuations and starts.\\
\((\lambda_g,\lambda_d)\) & Deterministic from \((\hat\theta_1,\hat\theta_0)\) via \eqref{eq:mstep_lambda} & Free: numerical, see \eqref{eq:Q_lambda_free}.\\
\bottomrule
\end{tabular}
\end{center}
All M-step updates maximise \(Q\): \(\alpha\), \(\pi\), and \(\sigma^2\) in closed form; \(\theta_1\) and \(\theta_0\) via a cheap one-dimensional Newton--Raphson (2--5 iterations) on the joint first-order conditions that collect gradient contributions from both the Markov prior and the EMG emission blocks.
\newpage






\end{document}